%% =====================================================================
%%  Corrigés — Unit 7 « Do you like coffee? »
%% =====================================================================
\corrigesUnite{7}{Do you like coffee?}

\begin{corrige}{1}
\begin{reponses}[4]
\item \rep{cheese}
\item \rep{water}
\item \rep{bread}
\item \rep{chicken}
\item \rep{eggs}
\item \rep{juice}
\item \rep{rice}
\item \rep{fish}
\end{reponses}
\piege{confondre \en{chicken} \phon{ˈtʃɪkɪn} (le poulet) et \en{kitchen} \phon{ˈkɪtʃɪn} (la cuisine,
la pièce).}
\end{corrige}

\begin{corrige}{2}
\begin{reponses}[1]
\item \rep{I don't like tea.}
\item \rep{She doesn't drink coffee.}
\item \rep{We don't eat meat.}
\item \rep{Jack doesn't cook every day.}
\item \rep{They don't play tennis.}
\item \rep{Emma doesn't like coffee.}
\end{reponses}
\pourquoi{avec \en{she}, \en{Jack}, \en{Emma}, on emploie \en{doesn't}, et le \en{-s} du verbe
disparaît : \en{drinks} → \en{doesn't drink}.}
\end{corrige}

\begin{corrige}{3}
\begin{reponses}[1]
\item \rep{Do you like jazz?}
\item \rep{Does she work here?}
\item \rep{Do they eat fish?}
\item \rep{Does Sam play the guitar?}
\item \rep{Does Jack cook?}
\item \rep{Do you drink tea?}
\end{reponses}
\retiens{question au présent simple : \en{Do} / \en{Does} + sujet + base verbale + reste de la phrase.}
\end{corrige}

\begin{corrige}{4}
\begin{reponses}[2]
\item \rep{Yes, she does.}
\item \rep{No, she doesn't.}
\item \rep{No, he doesn't.}
\item \rep{Yes, she does.}
\item \rep{No, she doesn't.}
\item \rep{Yes, he does.}
\end{reponses}
\piege{\en{Yes, she likes.} : la réponse courte reprend l'auxiliaire \en{does}, pas le verbe.}
\end{corrige}

\begin{corrige}{5}
\begin{reponses}[4]
\item \rep{Do}
\item \rep{do}
\item \rep{don't}
\item \rep{Does}
\item \rep{doesn't}
\item \rep{does}
\item \rep{doesn't}
\item \rep{don't}
\end{reponses}
\pourquoi{\en{Jack} et \en{Sam} → \en{does} / \en{doesn't} ; \en{you} et \en{I} → \en{do} / \en{don't}.
Ligne 6 : « il n'aime pas » → \en{doesn't like}.}
\end{corrige}

\begin{corrige}{6}
\begin{reponses}[3]
\item \rep{reading}
\item \rep{cooking}
\item \rep{dancing}
\item \rep{walking}
\item \rep{watching}
\item \rep{listening}
\end{reponses}
\pourquoi{\en{dance} perd son \en{e} : \en{dancing}. Les autres verbes ajoutent simplement \en{-ing}.}
\end{corrige}

\begin{corrige}{7}
\begin{reponses}[2]
\item \rep{Do you}
\item \rep{Does she}
\item \rep{What do you}
\item \rep{Do they}
\item \rep{Does he}
\item \rep{Where do you}
\end{reponses}
\retiens{à vitesse normale, \en{do you} se réduit à \phon{dəjə} et \en{does she} à \phon{dəʃi}.
L'accent tombe sur les mots porteurs de sens. Texte de la piste page~\pageref{trans-7.10}.}
\end{corrige}

\begin{corrige}{8}
\begin{reponses}[3]
\item \rep{no}
\item \rep{yes}
\item \rep{yes}
\item \rep{yes}
\item \rep{no}
\item \rep{no}
\item \rep{yes}
\item \rep{yes}
\item \rep{no}
\end{reponses}
\pourquoi{ligne par ligne : Emma (café, thé, cuisine), Jack, puis Sam. Sam est barista : il aime le
café, et aussi le thé (\en{I like tea too}) ; mais c'est sa sœur qui cuisine pour lui.}
\end{corrige}

\begin{corrige}{9}
\begin{reponses}[4]
\item \rep{Do}
\item \rep{do}
\item \rep{Does}
\item \rep{don't}
\item \rep{doesn't}
\item \rep{What}
\item \rep{cook}
\item \rep{make}
\end{reponses}
\retiens{\en{your family} (« ta famille ») est repris par \en{they} : \en{No, they don't.}}
\end{corrige}

\begin{corrige}{10}
\begin{reponses}[1]
\item \rep{Do you like coffee? — Yes, I do.}
\item \rep{Emma doesn't like coffee. She drinks tea.}
\item \rep{Does Jack cook? — No, he doesn't.}
\item \rep{I love reading and traveling.}
\item \rep{What do you do on weekends?}
\item \rep{We don't eat meat.}
\item \rep{What kind of music do you like?}
\item \rep{She hates cooking.}
\end{reponses}
\pourquoi{phrase 6 : « de viande » ne se traduit pas : \en{meat}, sans article. Phrase 4 : deux
verbes en \en{-ing} après \en{love}.}
\end{corrige}

\begin{corrige}{11}
Réponse modèle :
\begin{anglais}
\begin{tcolorbox}[otBase,colback=otGreenTint,colframe=otGreenTint,fontupper=\fontsize{10.4}{14}\selectfont]
I love coffee and I drink it every morning. I don't like tea. I like pasta and cheese, but I hate
fish! On weekends, I love reading and walking in the park. I don't like watching TV.
\end{tcolorbox}
\end{anglais}
\end{corrige}

\begin{corrige}{12}
Réponses modèles :
\begin{reponses}[1]
\item \rep{Yes, I do.} \textit{ou} \rep{No, I don't.}
\item \rep{I drink coffee.} \textit{ou} \rep{I drink tea and orange juice.}
\item \rep{Yes, I do. I cook pasta.} \textit{ou} \rep{No, I don't.}
\item \rep{I like jazz and pop.}
\item \rep{Yes, I'd love to!}
\end{reponses}
\end{corrige}

\begin{corrigetest}
\begin{reponses}[5][french]
\item \rep{b}
\item \rep{a}
\item \rep{c}
\item \rep{b}
\item \rep{a}
\item \rep{c}
\item \rep{a}
\item \rep{b}
\item \rep{c}
\item \rep{a}
\end{reponses}
\pourquoi{question 3 : après \en{doesn't}, le verbe reste à la base. Question 7 : pas d'article
pour parler du fromage en général.}
\end{corrigetest}
